\documentclass[11pt,a4paper]{article}
\usepackage{fontspec,amsmath,amssymb,amsthm,microtype}
\setmainfont{lmroman10-regular.otf}[BoldFont=lmroman10-bold.otf,ItalicFont=lmroman10-italic.otf,BoldItalicFont=lmroman10-bolditalic.otf]
\usepackage[margin=25mm]{geometry}
\usepackage[unicode,hidelinks]{hyperref}
\hypersetup{pdfauthor={LangevinLean Project}}
\setlength{\emergencystretch}{3em}
\setlength{\parskip}{0.4em}
\setlength{\parindent}{0pt}
\allowdisplaybreaks[2]
\raggedbottom
\title{\texorpdfstring{Construct the kinetic integral equation by Picard iteration}{Construct the kinetic integral equation by Picard iteration}}
\author{LangevinLean Project}
\date{}
\begin{document}\maketitle
Draft mathematical exposition; these notes may contain mistakes.


For existence, the phase product may use its convenient maximum norm. This norm is equivalent to the Euclidean phase norm but is used here only to establish a global Lipschitz vector field. The input noise is any continuous additive velocity path starting at zero. The theorem therefore constructs paths without postulating a reference kinetic process; a full Brownian solution theorem additionally identifies the stochastic noise and its laws.



The construction is pathwise and works for arbitrary continuous additive forcing. This provides the deterministic solution map on which the later Brownian law and coupling constructions are built.



Earlier results used here: \cite[Section~3.14]{LangevinLean2026ULASC}.

\section{\texorpdfstring{Write the drift in the shared integral-\allowbreak{}equation sign convention}{Write the drift in the shared integral-equation sign convention}}


The kinetic drift is \((-\gamma v-ug(x),v)\). The shared Picard theorem solves an equation with a minus sign in front of its drift integral, so the negative drift stores \((\gamma v+ug(x),-v)\). This sign choice ensures that \(Y_t-Y_0=-\int_0^tD(Y_s)ds+(B_t,0)\) has velocity acceleration \(-\gamma v-ug(x)\) and position derivative \(v\). The second noise coordinate is exactly zero.



Writing the equation componentwise removes any ambiguity: \(v_t=v_0-\int_0^t(\gamma v_s+ug(x_s))ds+B_t\) and \(x_t=x_0+\int_0^tv_sds\). In particular the position has an ordinary time derivative even when the forcing is only continuous.


\section{\texorpdfstring{Prove a global Lipschitz constant for phase space}{Prove a global Lipschitz constant for phase space}}


If \(g\) is \(L\)-Lipschitz and \(\gamma,u,L\ge0\), write \(R=\|(v,x)-(w,y)\|_{\max}\). Each coordinate difference is at most \(R\), and \(\|g(x)-g(y)\|\le LR\). The velocity component of the negative drift is bounded by \((\gamma+uL)R\) using the triangle inequality. The position component is bounded by \(R\). Hence the maximum is bounded by \((\gamma+uL+1)R\). The argument expands the product norm maximum, proves the two coordinate bounds, rewrites the drift difference into scalar multiples of differences with vector-space algebra, and combines the nonnegative coefficients. The constant is deliberately sufficient for existence rather than optimized for contraction.



The maximum norm is useful because each coordinate projection has operator norm one. A sharper phase Lipschitz constant would also prove existence, but its value is immaterial for the later contraction rate: that rate is derived separately from curvature in the transformed Hilbert norm.


\section{\texorpdfstring{Invoke the proved continuous-\allowbreak{}path Picard construction}{Invoke the proved continuous-path Picard construction}}


Let the underlying normed space be complete, \(B\) continuous with \(B_0=0\), and \(T\ge0\). The shared theorem the continuous-input Picard existence theorem constructs a continuous phase path \(Y\), with \(Y_0=\text{initial}\), satisfying the kinetic integral identity for every \(0\le t\le T\). The call supplies the concrete negative drift and its proved Lipschitz constant. It packages the noise as \((B_t,0)\), whose continuity follows from product continuity and whose value at zero follows from the zero initial forcing condition. This is finite-horizon pathwise existence; Gaussian stationarity, a canonical Brownian law, and full probabilistic chain convergence are not conclusions of this theorem.



Picard iteration solves the integral equation on a short interval by contraction of the integral map. Global Lipschitz continuity allows successive intervals to be joined without finite-time blowup. On a prescribed finite horizon the resulting path is unique and depends continuously on the initial phase and the forcing.



Earlier results used here: \cite[Section~3.14.3]{LangevinLean2026ULASC}.

\begin{thebibliography}{ULA-LSI}
\bibitem[ULA-SC]{LangevinLean2026ULASC}
LangevinLean Project. \emph{Unadjusted Langevin under Strong Convexity}. LL-01, 2026. Draft report, content revision 8d6335dce333. \url{https://odakyildiz.github.io/LangevinLean/docs/pdf/UnadjustedLangevin/FullProof.pdf}.
\end{thebibliography}

\end{document}
